% Esquisse d'un programme (cote 119, copy 1): archive cote 119, site page 50 (PDF page 51), typescript page 48.
% typescript. Draft compared: EsquisseFr.pdf (text layer):1725-1779.
% AI-transcribed by Claude Opus 5.5 (claude-opus-5-5, adaptive thinking, effort high), Message Batches API, batch esq1 (msgbatch_01LNsogFyF3y6KSp9UbtJtBk), 2026-11-12; not yet checked by eye.
\scanpage{119/50}
\mspage{48}

dant à un niveau très élémentaire. Dernièrement, j'ai finalement,
consacré un mois de réflexion intensive au développement d'une
construction purement combinatoire d'une sorte de ``surface modulaire''
associée à un système de $n$ pseudo-droites, qui classifie les
différentes ``positions relatives'' possibles (stables ou non) d'une
$(n+1)$ème pseudo-droite par rapport au système donné, ou encore :
les différentes ``affinisations'' possibles de ce système, par les
différents choix possibles d'une ``pseudo-droite à l'infini''. J'ai
l'impression d'avoir mis le doigt sur un objet remarquable, fai-%
sant apparaître un ordre imprévu dans des questions de classifi-%
cation qui jusqu'à présent apparaissaient assez chaotiques !
Mais ce n'est pas le lieu dans le présent rapport de m'étendre plus
à ce sujet.

Depuis 1977, dans toutes les questions (comme dans ces
deux derniers thèmes que je viens d'évoquer) où interviennent des
cartes bidimensionelles, la possibilité de les réaliser canoni-%
quement sur une surface conforme, donc sur une courbe algébrique
complexe dans le cas orienté compact, reste en filigrane constant
dans ma réflexion. Dans pratiquement tous les cas (en fait, tous
les cas sauf celui de certaines cartes sphériques avec ``peu d'auto-%
morphismes'') une telle réalisation conforme implique en fait une
\uline{métrique riemanienne canonique}, ou du moins, canonique à une cons-%
tante multiplicative près. Ces nouveaux éléments de structure
(sans même prendre en compte l'élément arithmétique, dont il a
été question au par. 3) sont de nature à transformer profondément
l'aspect initial des questions abordées, et les méthodes d'appro-%
che. Un début de familiarisation avec les belles idées de Thurston
sur la construction de l'espace de Teichmüller, en termes d'un jeu
très simple de chirurgie riemanienne hyperbolique, me confirme
dans ce pressentiment. Malheureusement, le niveau de culture très
modeste de presque tous les élèves qui ont travaillé avec moi pen-%
dans ces dix dernières années ne me permet pas d'aborder avec eux,
ne serait-ce que par allusion, de telles possibilités, alors que
l'assimilation d'un langage combinatoire minimum se heurte déjà,

\margin{50}
